\section{Conclusion}
\label{sec:conclusion}

For every item an agent remembers, three coupled questions---in what representation, on which tier, and
when to (re)materialize it---are today answered separately by three research communities. We gave a single
budget-constrained formulation that unifies them under one serve-time knob, and then asked whether the
resulting decision is learnable.

The answer is neither yes nor no but a decoupling, and that is the contribution. Reuse is predictable on
production traces and becomes more so with data, yet the placement benefit that predictability buys stays
near zero across the whole range; ordering, not calibration, is what the outcome turns on, and an oracle's
ordering would recover almost all of the gap. Reactive tiering, forecasting nothing, already delivers most
of the achievable benefit on one trace and rather less on the other, with the eviction signal itself
inverting between sources.

We arrived here only after removing ten measurement defects, each of which had
produced a plausible positive result; we report those mechanisms and the controls that catch them, because
the same traps are available to anyone measuring a memory scheduler. The useful next move is not a bigger model. It is to calibrate the cost
model against real tiers, to sweep the parameters an experiment fixes for convenience before believing any
result that depends on them, and to ask what a scheduler could do that ordering alone cannot. Agent memory
is ready to be scheduled---but first it has to be measured in a way that can tell schedulers apart.
